% Options for packages loaded elsewhere
\PassOptionsToPackage{unicode}{hyperref}
\PassOptionsToPackage{hyphens}{url}
\documentclass[
]{article}
\usepackage{xcolor}
\usepackage{amsmath,amssymb}
\setcounter{secnumdepth}{-\maxdimen} % remove section numbering
\usepackage{iftex}
\ifPDFTeX
  \usepackage[T1]{fontenc}
  \usepackage[utf8]{inputenc}
  \usepackage{textcomp} % provide euro and other symbols
\else % if luatex or xetex
  \usepackage{unicode-math} % this also loads fontspec
  \defaultfontfeatures{Scale=MatchLowercase}
  \defaultfontfeatures[\rmfamily]{Ligatures=TeX,Scale=1}
\fi
\usepackage{lmodern}
\ifPDFTeX\else
  % xetex/luatex font selection
\fi
% Use upquote if available, for straight quotes in verbatim environments
\IfFileExists{upquote.sty}{\usepackage{upquote}}{}
\IfFileExists{microtype.sty}{% use microtype if available
  \usepackage[]{microtype}
  \UseMicrotypeSet[protrusion]{basicmath} % disable protrusion for tt fonts
}{}
\makeatletter
\@ifundefined{KOMAClassName}{% if non-KOMA class
  \IfFileExists{parskip.sty}{%
    \usepackage{parskip}
  }{% else
    \setlength{\parindent}{0pt}
    \setlength{\parskip}{6pt plus 2pt minus 1pt}}
}{% if KOMA class
  \KOMAoptions{parskip=half}}
\makeatother
\setlength{\emergencystretch}{3em} % prevent overfull lines
\providecommand{\tightlist}{%
  \setlength{\itemsep}{0pt}\setlength{\parskip}{0pt}}
\usepackage{bookmark}
\IfFileExists{xurl.sty}{\usepackage{xurl}}{} % add URL line breaks if available
\urlstyle{same}
\hypersetup{
  hidelinks,
  pdfcreator={LaTeX via pandoc}}

\author{}
\date{}

\begin{document}

\section{Local specialization, cohomology, and the full lattice
spectrum}\label{local-specialization-cohomology-and-the-full-lattice-spectrum}

Date: 2026-09-23. This is an independent derivation of the local
statements in TS27--32, including TS28a, of
\texttt{TWO\_SOURCE\_SUPPORTED\_ZERO.tex}. That file was read in full
for its conventions; this task addresses only the specified local
analytic family, its cohomology and representation, and the map of full
supported-zero prime spectra.

The original source constants and metric entries are retained:

{[} A=-J\_2-\frac{1089}{4}C\textless0,\qquad
B=-48J\_2-4868C\textless0,\qquad
R=-\frac{\sqrt{\pi/32}}{32}e\^{}\{-8\}\textless0, {]}

{[}
50\textless K\_c,K\_s\textless64,\qquad G=\operatorname{diag}(K\_c,K\_s),
\qquad T\_*:=\frac{(R/A)^2}{4\pi}\textgreater0. \tag{LS1} {]}

Here (C,J\_2) are the original theta moments, with theta heat time fixed
at (1/32). The variable (T) belongs to the separate original
supported-integer moment action. The constant (R) in this note is the
original integral moment in TS12, not a source residual function. The
already certified inequalities in LS1 are inputs from the original
two-source calculation, rather than new numerical evaluations in this
note.

\subsection{1. The actual analytic disc and retained
coefficient}\label{the-actual-analytic-disc-and-retained-coefficient}

Use exactly

{[}
\Omega=\{T\in\mathbb C:\textbar T-T\_\emph{\textbar\textless T\_}/2\},
{]}

and the holomorphic square root positive at (T\_*). This branch exists
on the whole disc because (\Omega) lies in the right half-plane. Its
values have positive real part there, so neither (\sqrt T) nor
(\sqrt T+\sqrt{T_*}) vanishes.

Keep the original formulas

{[} c(T)=\frac{R(4\pi T)^{-1/2}-A}{B}, \qquad D(T)=K\_c+K\_sc(T)\^{}2,
{]}

{[} N(T)=\frac{K_sc(T)}{D(T)}

\begin{pmatrix}c(T)&-1\\c(T)^2&-c(T)\end{pmatrix}

, \quad u(T)=\binom{1}{c(T)},\quad w(T)=\binom{-K_sc(T)}{K_c}. \tag{LS2}
{]}

There is no denominator singularity on (\Omega). Indeed,

{[} 16(J\_2+1089C/4)\textless48J\_2+4868C {]}

because the difference of the right side and left side is exactly
(32J\_2+512C\textgreater0). Hence (0\textless A/B\textless1/16). Since
(R/\sqrt{4\pi T_*}=A),

{[} c(T)=\frac AB\left(\sqrt{\frac{T_*}{T}}-1\right), {]}

where the square root is the specified branch. For (T\in\Omega),
(\textbar T\textbar\textgreater T\_*/2), and therefore

{[}
\textbar c(T)\textbar\textless{}\frac{1+\sqrt2}{16}\textless{}\frac14.
{]}

Consequently (\textbar K\_sc(T)\^{}2\textbar\textless4\textless K\_c),
so
(\textbar D(T)\textbar{}\ge K\_c-\textbar K\_sc(T)\^{}2\textbar\textgreater46).
In particular

{[} \mathcal F(T)={[}u(T),w(T){]} =

\begin{pmatrix}1&-K_sc(T)\\c(T)&K_c\end{pmatrix}
\tag{LS3}

{]}

is holomorphic and invertible on the entire stated disc, with its exact
determinant (D(T)). This is the original frame, retaining the length of
each vector.

Put (\tau\emph{T=T-T}*), and retain the complete factorization

{[}

\begin{aligned}
c(T)&=-\frac{R\tau_T}
{B\sqrt{4\pi}\sqrt T\sqrt{T_*}(\sqrt T+\sqrt{T_*})},\\
n(T):=-K_sc(T)&=\tau_Tv(T),\\
v(T)&=\frac{K_sR}
{B\sqrt{4\pi}\sqrt T\sqrt{T_*}(\sqrt T+\sqrt{T_*})}.
\end{aligned}
\tag{LS4}

{]}

To derive LS4, subtract (R/(\sqrt{4\pi}\sqrt{T_*})=A) from
(R/(\sqrt{4\pi}\sqrt T)), and use
(\sqrt{T_*}-\sqrt T=-(T-T\_*)/(\sqrt T+\sqrt{T_*})). The function (v(T))
is holomorphic and nowhere zero on (\Omega). Its exact value at the
centre is

{[} v\_\emph{:=v(T\_})=n'(T\_*)
=\frac{K_sR}{2B\sqrt{4\pi}\,T_*^{3/2}}\textgreater0. \tag{LS5} {]}

No unit or sign in LS4--5 is omitted from the differential. Direct
multiplication gives

{[} N(T)u(T)=0,\qquad N(T)w(T)=n(T)u(T), \qquad
\mathcal F(T)\^{}\{-1\}N(T)\mathcal F(T) =

\begin{pmatrix}0&n(T)\\0&0\end{pmatrix}

. \tag{LS6} {]}

Thus (N(T)\^{}2=0) on the complete disc. At nonreal parameters these are
holomorphic algebraic identities. The continued formula for the
projection used to define (N) is not being declared Hermitian at a
nonreal parameter.

\subsection{2. The germ ring and local
cohomology}\label{the-germ-ring-and-local-cohomology}

Let (\mathcal O=\mathcal O\_\{\mathbb C,T\_\emph{\}) be the ring of
holomorphic germs at (T\_}), and let
(\mathfrak m=\tau\emph{T\mathcal O). Its residue field is
(k=\mathcal O/\mathfrak m\cong\mathbb C), by evaluation at (T}*).

The structural properties used here follow directly from convergent
Taylor expansions. A nonzero germ has a unique finite order of vanishing
(r\ge0), and is (\tau\emph{T\^{}ra(T)) with (a(T}*)\ne0). The latter
germ is a unit because its reciprocal is holomorphic on a sufficiently
small neighbourhood. Products of two nonzero germs have the sum of their
finite vanishing orders and a nonzero leading coefficient. Thus
(\mathcal O) is a domain. Its nonunits are precisely (\mathfrak m), so
it is local. Every nonzero ideal is generated by a germ of the smallest
vanishing order among its nonzero elements: division by the chosen
generator gives a holomorphic germ for every other element. In
particular (\mathcal O) is a discrete valuation ring in this explicit
one-variable sense.

By LS4 the actual coefficient (n=\tau\_Tv) is nonzero and satisfies

{[} (n)=(\tau\_T)=\mathfrak m, \tag{LS7} {]}

with the exact unit (v) still specified in LS4. The equality of ideals
does not change the matrix of the original differential in LS6.

Consider the two-periodic cochain complex, viewed in every integer
degree,

{[} \mathcal C\^{}j=\mathcal O\^{}2,
\qquad d\^{}j=N(T):\mathcal C\^{}j\longrightarrow\mathcal C\^{}\{j+1\}.
\tag{LS8} {]}

The chosen frame gives a unique expression (a u+b w) for every element
of every term. Formula LS6 gives (d\^{}j(a u+b w)=nbu). Since (n\ne0)
and (\mathcal O) is a domain,

{[} \ker d\^{}j=\mathcal O u,\qquad
\operatorname{im}d\^{}\{j-1\}=n\mathcal O u. {]}

Therefore in every integer degree there is the specified isomorphism

{[} \boxed{
H^j(\mathcal C^\bullet)=\mathcal O u/n\mathcal O u,
\qquad
\mathcal O/(n)\xrightarrow{\ \sim\ }H^j,
\quad [a]\longmapsto[a u].} \tag{LS9} {]}

Combining this actual coefficient quotient with LS7 gives
(H\^{}j\cong\mathcal O/\mathfrak m\cong\mathbb C), where the last map
evaluates the coefficient (a) at (T\_*). Its annihilator is exactly
((n)=\mathfrak m): an element kills the class of (u) if and only if it
belongs to that ideal.

The sheaf version on (\Omega) has the same explicit description, with a
sheaf supported only at (T\_*). Away from that point, (n) is invertible.
There is then an actual contracting homotopy, whose matrix in the
retained frame is

{[} h\^{}j=

\begin{pmatrix}0&0\\1/n(T)&0\end{pmatrix}

: \mathcal C\^{}j\longrightarrow\mathcal C\^{}\{j-1\}. {]}

Multiplying the matrices proves
(d\textsuperscript{\{j-1\}h}j+h\textsuperscript{\{j+1\}d}j=I). This
proves acyclicity on the punctured disc and gives the exact inverse
coefficient with all of LS4's factors retained.

\subsection{3. Fibre cohomology, Tor, and the retained Bockstein
factor}\label{fibre-cohomology-tor-and-the-retained-bockstein-factor}

At (T=T\_*) the two frame vectors are exactly

{[} u\_\emph{=(1,0)\^{}\{\mathsf T\},\qquad
w\_}=(0,K\_c)\^{}\{\mathsf T\}. \tag{LS10} {]}

Since (n(T\_*)=0), the evaluated differential is zero. Thus the
evaluated complex has

{[} \boxed{H^j(\mathcal C^\bullet\otimes_{\mathcal O}k)
=k u_*\oplus k w_*\cong\mathbb C^2} \tag{LS11} {]}

in every degree. In contrast, LS9 gives
(H\^{}j(\mathcal C\^{}\bullet)\otimes k\cong k). The natural map from
that tensor product to LS11 sends the coefficient class ({[}a{]}) to
(a(T\_\emph{)u\_}). Its image is precisely the first line in LS11.

The additional class is measured by the actual free resolution

{[} 0\longrightarrow\mathcal O r\_1 \xrightarrow{\ n\ }\mathcal O r\_0
\longrightarrow\mathcal O/(n)\longrightarrow0. \tag{LS12} {]}

Injectivity of its first arrow follows from the domain property. After
tensoring with (k), multiplication by (n) becomes zero, and hence

{[} \operatorname{Tor}\^{}\{\mathcal O\}\_1(\mathcal O/(n),k)=k r\_1.
{]}

Using this resolution for (H\^{}\{j+1\}), the entire base-change
sequence, with its maps, is

{[} \boxed{
0\longrightarrow k[u]\xrightarrow{\ i\ }
k u_*\oplus k w_*\xrightarrow{\ q\ }k r_1
\longrightarrow0,
\quad i(a)=a u_*,\quad q(a u_*+b w_*)=b r_1.} \tag{LS13} {]}

The compatibility with the differential is exactly (d(w)=n u), so the
displayed quotient is the Tor term from LS12. Its kernel and
surjectivity can also be checked directly in the two-dimensional vector
space. The chosen frame supplies the explicit splitting
(r\_1\mapsto w\_*); no assertion of a frame-independent preferred
splitting is needed.

There is a second useful exact description that retains the unit in LS4
visibly. The short exact sequence of complexes

{[} 0\longrightarrow\mathcal C\^{}\bullet
\xrightarrow{\ \tau_T\ }\mathcal C\^{}\bullet
\longrightarrow\mathcal C\^{}\bullet\otimes k\longrightarrow0 {]}

has connecting map
(\delta:H\^{}j(\mathcal C\^{}\bullet\otimes k)\to H\^{}\{j+1\}(\mathcal C\^{}\bullet)).
Lift (a u\_\emph{+b w\_}) to (a u+b w). Its differential is (bn
u=\tau\_T bv u). Division by the exact displayed (\tau\_T), followed by
taking its class, gives

{[} \boxed{\delta(a u_*+b w_*)=b v_*[u],
\qquad v_*=
\frac{K_sR}{2B\sqrt{4\pi}\,T_*^{3/2}}>0.} \tag{LS14} {]}

Thus when the right side of LS13 is identified with (H\^{}\{j+1\})
through the parameter-based connecting map, its coefficient is
(v\_\emph{). The quotient map using the exact (n)-resolution in LS12 has
coefficient one by its stated choice of generator. These are two proved
coordinate descriptions of the same one-dimensional additional
contribution, related by multiplication by the retained nonzero number
(v\_}).

This verifies TS28 and TS28a, including the extra fibre class. The fibre
cohomology has dimension two and is not identified with the
one-dimensional tensor product of the germ cohomology.

\subsection{4. The image algebra, its minors, and its exact
specialization}\label{the-image-algebra-its-minors-and-its-exact-specialization}

Define the original representation

{[} \rho:\mathcal A:=\mathcal O{[}\eta{]}/(\eta\^{}2)
\longrightarrow\operatorname{End}\_\{\mathcal O\}(\mathcal O\^{}2),
\qquad \rho(\eta)=N(T). \tag{LS15} {]}

In the frame (\mathcal F), the full matrix of (\rho(a+b\eta)) is

{[}

\begin{pmatrix}a&bn\\0&a\end{pmatrix}

. \tag{LS16} {]}

The map is multiplicative because (N\^{}2=0), and it is injective: a
zero matrix in LS16 has (a=0), then (bn=0), and finally (b=0). Thus its
image (\mathcal A\_\{\rm im\}) is a free rank-two (\mathcal O)-algebra
with the explicit basis (I,N). Both the abstract algebra and this image
algebra retain the original relation (\eta\^{}2=0).

To inspect the embedding as a map of (\mathcal O)-modules, order the
four entries as ((11,12,21,22)). The exact two-column matrix is

{[} M\_\rho=

\begin{pmatrix}1&0\\0&n\\0&0\\1&0\end{pmatrix}

. \tag{LS17} {]}

Subtracting the first row from the fourth is an invertible row
operation; the two nonzero remaining diagonal entries are (1,n). The six
two-row minors, in the order ((12),(13),(14),(23),(24),(34)), are
exactly

{[} n,\quad0,\quad0,\quad0,\quad-n,\quad0. \tag{LS18} {]}

Consequently their ideal is ((n)=\mathfrak m), with its actual generator
(n=\tau\_Tv). The minor (n) has a simple zero with the exact positive
derivative LS5. The one-entry minors generate the unit ideal because the
matrix contains the entry one. These determinantal ideals are unchanged
by the invertible original frame change.

The module cokernel is completely explicit. Define

{[} \Psi:\operatorname{Mat}\_2(\mathcal O)
\longrightarrow\mathcal O\oplus\mathcal O\oplus\mathcal O/(n), \quad

\begin{pmatrix}p&q\\r&s\end{pmatrix}

\longmapsto(r,s-p,{[}q{]}). \tag{LS19} {]}

It is surjective: choose (p=0), then choose the indicated (r,s) and any
representative (q). Its kernel consists exactly of (r=0), (s=p), and
(q=bn), which is LS16. Therefore

{[} \boxed{\operatorname{coker}_{\mathcal O}\rho
\cong\mathcal O^2\oplus\mathcal O/(n).} \tag{LS20} {]}

This is the cokernel of the underlying module embedding, not an algebra
quotient by a ring ideal. The image is a unital subalgebra, and contains
the identity, so interpreting that module cokernel as a quotient algebra
would be incorrect.

The abstract algebra fibre remains

{[} \mathcal A\otimes\_\{\mathcal O\}k \cong k{[}\eta{]}/(\eta\^{}2),
{]}

and the free image algebra has the identical abstract fibre
(\mathcal A\_\{\rm im\}/\mathfrak m\mathcal A\_\{\rm im\}). In the
latter, the class of (N) is nonzero. Indeed if (N=\tau\_T(aI+bN)),
comparing its diagonal entries gives (a=0), and a nonzero entry of (N)
then gives (1=\tau\_Tb), impossible in the local ring.

But evaluating the embedding in the ambient matrices gives

{[}
\boxed{\rho_*:k[\eta]/(\eta^2)\longrightarrow\operatorname{Mat}_2(k),
\qquad a+b\eta\longmapsto aI,
\quad\ker\rho_*=(\eta),\quad\operatorname{im}\rho_*=kI.} \tag{LS21} {]}

Equivalently, inside the ambient matrix module
(\mathcal B=\operatorname{Mat}\_2(\mathcal O)),

{[} \mathcal A\_\{\rm im\}\cap\mathfrak m\mathcal B =\mathfrak m
I\oplus\mathcal O N, \qquad \mathfrak m\mathcal A\_\{\rm im\}
=\mathfrak m I\oplus\mathfrak m N. \tag{LS22} {]}

The first equality follows directly from LS16, because
(n\in\mathfrak m); the diagonal coefficient must lie in (\mathfrak m),
while the off-diagonal coefficient (bn) already does. Their quotient is
exactly the one-dimensional kernel in LS21.

Tensoring the exact embedding and its module cokernel with (k) records
this failure of injectivity as

{[} 0\longrightarrow k\eta \longrightarrow k{[}\eta{]}/(\eta\^{}2)
\xrightarrow{\ \rho_*\ }\operatorname{Mat}\_2(k)
\longrightarrow k\^{}3\longrightarrow0. \tag{LS23} {]}

The initial line is
(\operatorname{Tor}\_1\^{}\{\mathcal O\}(\mathcal O\^{}2\oplus\mathcal O/(n),k)),
calculated using LS12. The final space has the three coordinates in LS19
after reduction. The middle image has dimension one, so the complete
sequence is also verified directly. This proves the image-algebra,
minors, cokernel, and representation-specialization assertions of TS29
with their exact domains and maps.

\subsection{5. Prime classification with the complete lattice
retained}\label{prime-classification-with-the-complete-lattice-retained}

Let (L) be the original bounded distributive lattice with (0\_L\ne1\_L).
No quotient of (L) is taken. For a commutative ring (R\_0), define

{[} S\_L(R\_0)=G\_L(R\_0) =\{(r,1\_L):r\in R\_0\}
\cup\{z\_\lambda=(0,\lambda):\lambda\ne1\_L\}. \tag{LS24} {]}

The notation (z\_\{1\_L\}=e=(0,1\_L)) also denotes the top-supported
zero, and the additive identity is the distinct element
(\tau=z\_\{0\_L\}). For elements with their original coordinates,

{[} (r,\lambda)+(s,\mu)=(r+s,\lambda\vee\mu), \qquad
(r,\lambda)(s,\mu)=(rs,\lambda\wedge\mu). \tag{LS25} {]}

The carrier is closed under these operations: a nonzero amplitude in a
sum forces a top join, and a nonzero amplitude in a product forces both
input labels to be top. Associativity and the distributive laws follow
from the corresponding ring and distributive-lattice laws in the
displayed coordinates. A cancellation in the ring amplitude at top label
gives (e), retaining that label.

The amplitude map (\operatorname{amp}:S\_L(R\_0)\to R\_0) and the label
map (\operatorname{lab}:S\_L(R\_0)\to L) preserve these operations and
their identities, where the operations in (L) are join and meet. The
notation ({[}r{]}=(r,1\_L)) below is only an element notation: the map
(r\mapsto[r]) sends ring zero to (e), not to the semiring zero (\tau).

A semiring ideal here contains (\tau), is closed under addition, and
absorbs multiplication by every semiring element. A prime ideal is a
proper ideal (P) for which (ab\in P) implies (a\in P) or (b\in P).

For a ring prime (\mathfrak p\subset R\_0), put

{[}
P\_\{\mathfrak p\}\textsuperscript{\{R\_0\}=\operatorname{amp}}\{-1\}(\mathfrak p).
{]}

For a proper lattice ideal (J\subset L) prime for meets, put

{[} Q\_J\^{}\{R\_0\}=\{z\_\lambda:\lambda\in J\}
=\operatorname{lab}\^{}\{-1\}(J). \tag{LS26} {]}

A lattice ideal means a nonempty down-set closed under finite joins;
properness implies (1\_L\notin J). It is prime for meets when
(\lambda\wedge\mu\in J) forces (\lambda\in J) or (\mu\in J). Both types
in LS26 are prime by their defining maps. In particular the converse for
(Q\_J) uses the label meet and works even if the coefficient ring has
zero divisors; a top-supported product remains top-supported if its
amplitude vanishes.

These are all primes, with no lattice points identified. To prove it,
let (P) be a prime ideal excluding (e). It contains no top-supported
element ({[}r{]}): otherwise it contains ({[}-1{]}{[}r{]}={[}-r{]}),
hence also their sum (e). Thus (P) consists only of lower labelled
zeros. Its labels form a lattice ideal because addition takes joins, and
for (\mu\le\lambda), (z\_\mu=z\_\lambda z\_\mu). Primality of (P) is
exactly the meet-prime property. This gives (P=Q\_J\^{}\{R\_0\}).

If instead (e\in P), then (ez\_\lambda=z\_\lambda) puts every labelled
zero in (P). The set (\mathfrak p=\{r\in R\_0:{[}r{]}\in P\}) is a
proper ring ideal: top-supported sums give additive closure,
multiplication by ({[}-1{]}) gives negatives, and multiplication by
every ({[}r{]}) gives absorption. The properness and prime-product
property are inherited from (P). Thus
(P=P\_\{\mathfrak p\}\^{}\{R\_0\}). This proves the full classification

{[} \boxed{\operatorname{Spec}G_L(R_0)
=\{Q_J^{R_0}:J\in\operatorname{Spec}_{\rm lat}L\}
\ \sqcup\ \{P_{\mathfrak p}^{R_0}:\mathfrak p\in\operatorname{Spec}R_0\}.}
\tag{LS27} {]}

For a domain (R\_0), the original supported-zero prime is

{[} P\_e\textsuperscript{\{R\_0\}=P\_\{(0)\}}\{R\_0\}
=\{z\_\lambda:\lambda\in L\}=(e). \tag{LS28} {]}

The principal ideal equality follows from (ez\_\lambda=z\_\lambda),
while every product by (e) and every sum of such products has zero
amplitude. Domain primality follows either from LS27 or from the
original fact that a product has zero amplitude exactly when one
amplitude is zero.

The inclusions retain the complete lattice order:
(Q\_J\^{}\{R\_0\}\subseteq Q\_\{J'\}\^{}\{R\_0\}) holds exactly when
(J\subseteq J'). Every (Q\_J\^{}\{R\_0\}) is strictly below every
(P\_\{\mathfrak p\}\^{}\{R\_0\}), because the former omits (e) and the
latter contains all labelled zeros. Among the second type, inclusions
are precisely inclusions of ring primes. Thus distinct and incomparable
lattice primes remain distinct and incomparable.

The germ ring has only two ring primes, ((0)) and (\mathfrak m). It is a
domain, so ((0)) is prime. If a nonzero prime contains a nonzero germ
(\tau\_T\^{}ru), it cannot contain its unit (u). Primality then forces
(\tau\_T) into the prime, so the prime contains the maximal ideal
(\mathfrak m) and equals it. Hence the entire spectrum over the germ
ring is

{[} \{Q\_J\^{}\{\mathcal O\}:J\in\operatorname{Spec}\emph{\{\rm lat\}L\}
~\sqcup~\{P\_e\^{}\{\mathcal O\},P}\{\mathfrak m\}\^{}\{\mathcal O\}\},
\qquad Q\_J\^{}\{\mathcal O\}\subsetneq P\_e\^{}\{\mathcal O\}
\subsetneq P\_\{\mathfrak m\}\^{}\{\mathcal O\}. \tag{LS29} {]}

The last inclusion is strict because ({[}\tau\_T{]}) has nonzero germ
amplitude and belongs to (P\_\{\mathfrak m\}\^{}\{\mathcal O\}). If (L)
has exactly two elements, its sole lattice prime is (J=\{0\_L\}), giving
(Q\_J=\{\tau\}). For a general lattice, LS29 keeps every actual (J); no
single lower point replaces that family.

\subsection{6. The full prime-spectrum pullback under integer
constants}\label{the-full-prime-spectrum-pullback-under-integer-constants}

The ordinary unital ring map (\mathbb Z\to\mathcal O) sends an integer
to its constant holomorphic germ. It induces the exact label-preserving
unital semiring map

{[} \varphi:G\_L(\mathbb Z)\longrightarrow G\_L(\mathcal O),
\qquad [n]\longmapsto[n],\quad z\_\lambda\longmapsto z\_\lambda.
\tag{LS30} {]}

The common value for ({[}0{]}=e=z\_\{1\_L\}) is consistent. Addition and
multiplication are preserved because the amplitude map preserves ring
operations and every label is fixed.

The induced map on spectra is contraction
(P\mapsto\varphi\^{}\{-1\}(P)). Its complete values are

{[} \boxed{
(\operatorname{Spec}\varphi)(Q_J^{\mathcal O})=Q_J^{\mathbb Z},
\quad
(\operatorname{Spec}\varphi)(P_e^{\mathcal O})=P_e^{\mathbb Z},
\quad
(\operatorname{Spec}\varphi)(P_{\mathfrak m}^{\mathcal O})=P_e^{\mathbb Z}.}
\tag{LS31} {]}

For the first equality, membership is determined by the unchanged label
in the actual set (J). For the second, a constant germ is zero if and
only if its integer is zero. For the third, a constant germ vanishes at
(T\_*) if and only if its integer is zero. All lower zero labels occur
in both contractions of the latter two types. This proves the three
assertions directly as equalities of ideals.

Every nonzero integer (n) is a unit in (\mathcal O), with the actual
constant inverse (1/n). Its top-supported element ({[}n{]}) is
correspondingly a unit in (G\_L(\mathcal O)). Thus it lies in no proper
prime there. In particular no ordinary positive integer prime (p) can
appear as a ring-prime contraction ((p)), and no point
(P\_\{(p)\}\^{}\{\mathbb Z\}) lies in the image of LS31.

There are no additional prime points to consider: the ring primes of
(\mathbb Z) are ((0)) and ((p)) for positive prime integers (p). Indeed
a nonzero proper ideal has a least positive element (d), the division
algorithm makes the ideal ((d)), and primality forces (d) to be a prime
integer. Therefore the exact set-theoretic image of this germ chart is

{[} \boxed{\operatorname{im}(\operatorname{Spec}\varphi)
=\{Q_J^{\mathbb Z}:J\in\operatorname{Spec}_{\rm lat}L\}
\cup\{P_e^{\mathbb Z}\}.} \tag{LS32} {]}

The fibre over (P\_e\^{}\{\mathbb Z\}) consists of the two distinct
points (P\_e\^{}\{\mathcal O\}) and
(P\_\{\mathfrak m\}\^{}\{\mathcal O\}). The fibre over any
(Q\_J\^{}\{\mathbb Z\}) is its single identically labelled
(Q\_J\^{}\{\mathcal O\}); every ordinary integer-prime fibre is empty.

With the usual prime-ideal topology (V(I)=\{P:I\subseteq P\}), this map
is continuous. Directly, a prime (P) lies over a member of (V(I)) if and
only if (\varphi(I)\subseteq P), giving the inverse-image closed set
defined by the ideal generated by (\varphi(I)). This topology also
records that

{[} \overline{\{P_e^{\mathbb Z}\}} =V(P\_e\^{}\{\mathbb Z\})
=\{P\_e\^{}\{\mathbb Z\}\}\cup\{P\_\{(p)\}\^{}\{\mathbb Z\}:p\text{ a positive prime integer}\}.
\tag{LS33} {]}

Thus the ordinary integer-prime points are absent from the chart image,
but are in the closure of its supported-zero image point. No
lattice-prime point (Q\_J) belongs to this closure because it does not
contain (e).

\subsection{7. The cohomology support and its exact arithmetic
location}\label{the-cohomology-support-and-its-exact-arithmetic-location}

Use the germ cohomology module (H\^{}j=\mathcal O/(n)) with the
specified generator class ({[}u{]}). Regard it as a semimodule over
(S=G\_L(\mathcal O)) by the exact amplitude action

{[} s\cdot[a]={[}\operatorname{amp}(s)a{]}. \tag{LS34} {]}

Its additive structure remains that of the abelian group
(\mathcal O/(n)). Every labelled zero acts as zero. The annihilator is
evaluated exactly:

{[} \boxed{\operatorname{Ann}_{G_L(\mathcal O)}H^j
=\operatorname{amp}^{-1}((n))
=P_{\mathfrak m}^{\mathcal O}.} \tag{LS35} {]}

Indeed it is enough to test a scalar on the generator class ({[}u{]}),
and it kills that class precisely when its amplitude belongs to
((n)=\mathfrak m).

Here is a direct localization proof of the support, so no general
assertion about supports of semimodules is being assumed. At a prime
(Q\_J\^{}\{\mathcal O\}), the scalar (e) is outside the prime and acts
by zero. It becomes invertible in the localization, forcing the
localized module to be zero. At (P\_e\^{}\{\mathcal O\}), the element
({[}\tau\_T{]}) lies outside the prime and annihilates (H\^{}j), so this
localization is also zero. At (P\_\{\mathfrak m\}\^{}\{\mathcal O\}),
every scalar outside the prime has a unit germ amplitude. These scalars
already act invertibly on (\mathcal O/(n)), and localization leaves the
nonzero module unchanged. The last identification sends ({[}a{]}/s) to
({[}\operatorname{amp}(s)\^{}\{-1\}a{]}).

Consequently

{[} \boxed{\operatorname{Supp}_{G_L(\mathcal O)}H^j
=\{P_{\mathfrak m}^{\mathcal O}\},
\qquad
(\operatorname{Spec}\varphi)\bigl(\operatorname{Supp}_{G_L(\mathcal O)}H^j\bigr)
=\{P_e^{\mathbb Z}\}.} \tag{LS36} {]}

This proves the precise TS31--32 claim: the local cohomology support
point lies over the original supported-zero prime, every lattice prime
is fixed by the chart map, and ordinary integer primes are absent from
that map's image because their amplitudes become units.

The amplitude action is also forced by the additive-group type of this
cohomology. In any semimodule over (G\_L(R\_0)) whose addition is a
group, the equality (z\_\lambda+z\_\lambda=z\_\lambda) gives
(z\_\lambda m+z\_\lambda m=z\_\lambda m), so cancellation gives
(z\_\lambda m=0) for every (m). Thus every labelled zero acts as zero.
Defining (r m={[}r{]}m) then gives an ordinary (R\_0)-module action;
addition, multiplication and identity follow from the top-supported
operations, while ring zero acts as zero because ({[}0{]}=e). Conversely
an (R\_0)-module gives a semimodule by amplitude. These constructions
are inverse on both objects and homomorphisms.

Accordingly the nonzero original operator (E\_T) representing
multiplication by (e) on coefficient data is a representation of that
multiplicative action, not a semiring-linear scalar action on an
additive group: it cannot satisfy the additive relation (e+e=e) as a
nonzero endomorphism of a complex vector space. This establishes the
exact type distinction already stated in TS32 and preserves the original
operator rather than equating its action with the amplitude action.

\subsection{8. Restricting the module to integers is a different support
calculation}\label{restricting-the-module-to-integers-is-a-different-support-calculation}

The image assertion LS36 must not be strengthened into absence of
ordinary integer primes from the support of the module after restricting
scalars. The exact restricted module is

{[} H\^{}j\cong\mathbb C {]}

with integer scalars acting as ordinary integer multiplication and every
labelled zero acting as zero. Its annihilator over (G\_L(\mathbb Z)) is
(P\_e\^{}\{\mathbb Z\}), since a nonzero integer does not annihilate
(\mathbb C).

At each lattice prime (Q\_J\^{}\{\mathbb Z\}), localization again
inverts (e), which annihilates the module, so the localized module is
zero. At (P\_e\^{}\{\mathbb Z\}), all scalars outside the prime have
nonzero integer amplitudes, and each acts invertibly on (\mathbb C). The
localization is therefore (\mathbb C), not zero. At an ordinary integer
prime (P\_\{(p)\}\^{}\{\mathbb Z\}), each scalar outside the prime has
integer amplitude not divisible by (p), hence nonzero, and again acts
invertibly on (\mathbb C). This localization is likewise (\mathbb C).

It follows by direct computation that

{[} \boxed{
\operatorname{Supp}_{G_L(\mathbb Z)}(\operatorname{Res}_{\varphi}H^j)
=\{P_e^{\mathbb Z}\}\cup
\{P_{(p)}^{\mathbb Z}:p\text{ a positive prime integer}\}
=V(P_e^{\mathbb Z}).} \tag{LS37} {]}

Thus LS36 is the image of the original germ support, whereas LS37 is the
support of the restricted module. They have the exact relation

{[}
\operatorname{Supp}\emph{\{G\_L(\mathbb Z)\}(\operatorname{Res}}\{\varphi\}H\^{}j)
=\overline{(\operatorname{Spec}\varphi)
(\operatorname{Supp}_{G_L(\mathcal O)}H^j)} {]}

in this actual calculation, by LS33, LS36 and LS37. The closure is
strict because every ordinary positive integer prime gives an additional
point. This does not alter TS32's stated chart image; it specifies the
further map and its support behaviour exactly.

\subsection{Verification and provenance
log}\label{verification-and-provenance-log}

\begin{enumerate}
\def\labelenumi{\arabic{enumi}.}
\tightlist
\item
  Read \texttt{TWO\_SOURCE\_SUPPORTED\_ZERO.tex}, with the bounded task
  focused on TS27--32 and TS28a.
\item
  Proved nonvanishing of every denominator on the actual disc, retaining
  the original constants and square-root branch.
\item
  Computed the germ complex, its contracting homotopy off the centre,
  its cohomology, its fibre cohomology, and both the exact
  (n)-resolution and the parameter-based Bockstein with its full unit
  factor.
\item
  Computed the representation's image algebra, all two-column minors,
  its module cokernel, abstract algebra fibre, evaluated image, and the
  specialization kernel.
\item
  Proved the complete prime classification for the unchanged lattice and
  the full contraction map from integer constants.
\item
  Computed the cohomology support by localization and proved that it
  maps to (P\_e\^{}\{\mathbb Z\}).
\item
  Identified and proved the necessary additional distinction:
  restriction of the module to the integers has nonzero localizations at
  ordinary integer-prime points, although those points are absent from
  the chart image.
\end{enumerate}

No numerical certificate, whole-programme review, rendering, indexing,
coordination with other tasks, artifact packaging, or publication was
performed in this derivation.

\end{document}
